\documentclass[10pt]{article}

\usepackage[margin=0.8in]{geometry}
\usepackage{amsmath, amssymb}
\usepackage{hyperref}

\title{Ambiguity Resolution in Large Language Models:\\
A Decision-Theoretic Study of Pragmatic Inference and Clarification}

\author{
Thomas Lee, Ziyu Ge, Shizhe Zhang \\
\texttt{\{tholee, ziyuge, shizhe\_zhang\}@berkeley.edu}
}

\date{}

\begin{document}
\maketitle

\noindent
\textbf{Problem.}
Large language models (LLMs) frequently encounter ambiguous instructions, yet it remains unclear how they balance implicit inference against explicit clarification. We study ambiguity resolution in a controlled synthetic reference game.

Each instance consists of:
\begin{itemize}
    \item \textbf{Input:} A scene containing $N$ objects defined by structured attributes (color, shape, size, position), and a natural language description referring to a target object.
    \item \textbf{Output:} A predicted target object. The model may optionally ask at most one binary clarification question before committing.
\end{itemize}

Ambiguity is systematically controlled via attribute overlap and referential entropy $H(T \mid u)$, where $T$ denotes the target and $u$ the utterance.

Primary success metric is accuracy:
\[
\text{Accuracy} = \frac{\text{Correct Predictions}}{\text{Total Instances}}.
\]
Secondary metrics include expected utility under clarification cost $c$, Brier score, Expected Calibration Error (ECE), and clarification frequency.

\vspace{0.5em}
\noindent
\textbf{Method.}
We compare two decision mechanisms:

\textit{(1) Pragmatic Inference (LLM-RSA).}  
We approximate recursive reasoning inspired by Rational Speech Act (RSA) modeling:
\[
P_{L_1}(t \mid u) \propto P_{S_0}(u \mid t) P(t),
\]
where $P_{S_0}(u \mid t)$ is estimated using LLM likelihoods. The model selects the object with maximal posterior probability.

\textit{(2) Clarification-Augmented Strategy.}  
The model may ask one binary question before predicting. A decision rule compares expected utilities:
\[
\text{Ask if } \mathbb{E}[U_{\text{ask}}] > \mathbb{E}[U_{\text{commit}}],
\]
where utility incorporates both prediction accuracy and clarification cost $c$.

All experiments use fully synthetic scenes and prompt-based inference without fine-tuning, isolating reasoning behavior from memorized knowledge.

\vspace{0.5em}
\noindent
\textbf{Anticipated Results.}
We expect that:
\begin{itemize}
    \item Pure pragmatic inference performs well under low ambiguity but degrades as distractors increase.
    \item Clarification improves accuracy in high-entropy regimes when $c$ is moderate.
    \item LLMs may under- or over-ask relative to the decision-theoretic optimum, revealing miscalibration.
    \item Calibration metrics will expose whether confidence tracks ambiguity.
\end{itemize}

\vspace{0.5em}
\noindent
\textbf{Why It Matters.}
Ambiguity is central to real-world human--AI interaction. Systems that overcommit risk errors; systems that over-ask harm usability. Our framework provides a principled, decision-theoretic lens for analyzing when LLMs infer versus inquire. By isolating ambiguity under controlled conditions, this work contributes to understanding uncertainty handling, communicative rationality, and robustness in modern language models.

\end{document}